\section{Các hướng tiếp cận hiện nay}

\subsection{Hệ thống gia sư thông minh truyền thống}
Hệ thống gia sư thông minh truyền thống tách kiến thức chuyên môn, mô hình người học và chiến lược dạy thành các thành phần riêng. AutoTutor là một ví dụ tiêu biểu: hệ thống dẫn dắt học sinh qua hội thoại bằng ngôn ngữ tự nhiên theo kịch bản soạn trước~\cite{graesser2005}. Truy vết tri thức theo mô hình Bayes ước lượng xác suất học sinh đã nắm một kỹ năng từ chuỗi câu trả lời~\cite{corbett1995}, và các mô hình học sâu về sau mở rộng cách tiếp cận đó~\cite{piech2015}. Ưu điểm của dòng này là hành vi sư phạm kiểm soát được và mức hiểu bài có cơ sở tính toán rõ ràng. Nhược điểm nằm ở chi phí soạn nội dung và kịch bản hội thoại, vốn cứng và khó mở rộng ra những câu hỏi ngoài kịch bản.

\subsection{Gia sư dựa trên mô hình ngôn ngữ lớn}
Các công trình gần đây dùng mô hình ngôn ngữ để làm hội thoại dạy học linh hoạt hơn. MathDial xây dựng tập hội thoại dạy Toán có đặc tính sư phạm để huấn luyện và đánh giá gia sư~\cite{macina2023}. SocraticLM tinh chỉnh mô hình để dạy theo phương pháp Socratic~\cite{liu2024socratic}. Chowdhury và cộng sự kết hợp mô hình ngôn ngữ với khung hội thoại của AutoTutor để giữ được hàng rào sư phạm~\cite{chowdhury2024}. Các nhà cung cấp lớn cũng đưa ra những chế độ học riêng. Study mode của ChatGPT được mô tả là một bộ chỉ dẫn hệ thống tùy biến chứ không phải một mô hình riêng~\cite{studymode}. Guided Learning của Gemini dựa trên họ mô hình LearnLM~\cite{guidedlearning,jurenka2024}. Khanmigo gắn GPT-4 với kho bài của Khan Academy~\cite{khanmigo}.

\subsection{Khoảng trống nghiên cứu}
Các cách làm trên có chung một đặc điểm: mô hình ngôn ngữ vẫn là nơi đưa ra quyết định, còn ràng buộc sư phạm chỉ tồn tại dưới dạng lời khuyên trong ngữ cảnh hoặc dưới dạng trọng số đã tinh chỉnh. Trường hợp lời khuyên bị bỏ qua, chẳng hạn khi trẻ khăng khăng một đáp án sai hoặc nhờ làm hộ, không có cơ chế nào bảo đảm hành vi đúng. Ba khoảng trống cụ thể xuất hiện từ đó.

Khoảng trống thứ nhất liên quan tới kiểm soát: chưa có cơ chế tách quyết định sư phạm (hỏi gì, gợi ý ở mức nào, khi nào chuyển bài) khỏi quá trình sinh ngôn ngữ, nên khó chứng minh hành vi bất biến như không lộ đáp án. Khoảng trống thứ hai liên quan tới nền tảng nội dung: kiến thức của gia sư thường là kiến thức chung của mô hình, không bám tài liệu của giáo viên và không biết ranh giới giữa bài đang học với bài sắp học. Khoảng trống thứ ba liên quan tới đánh giá: mức hiểu bài thường được suy ra từ lời mô hình, khó tái lập và khó kiểm toán. Đồ án nhắm tới việc lấp ba khoảng trống này trong phạm vi hẹp của hai chương Toán lớp 5.
